% 攻读硕士学位期间取得的科研成果（与参考文献、致谢同级目录）
\newpage
\markboth{攻读硕士学位期间取得的科研成果}{攻读硕士学位期间取得的科研成果}
\phantomsection
\addcontentsline{toc}{section}{攻读硕士学位期间取得的科研成果}

\begin{center}
{\fontsize{15.75pt}{\baselineskip}\selectfont\heiti 攻读硕士学位期间取得的科研成果}
\end{center}
\vspace{0.35\baselineskip}

% 与中文摘要、正文一致：小四号宋体，行距 12pt/18pt（约 1.5 倍），西文与数字随正文主字体（Times New Roman）
\setlength{\parindent}{2em}
\setlength{\parskip}{0pt}
\begingroup
\fontsize{12pt}{18pt}\selectfont

% 不用 \subsubsection*，避免 titlesec 块级标题带来大块竖向留白
\noindent\heiti 参与科研项目及科研获奖\par
\vspace{0.35\baselineskip}
\setlength{\parindent}{0pt}
\songti
\begin{enumerate}[
  label={[\arabic*]},
  align=left,
  leftmargin=*,
  itemindent=0pt,
  labelindent=0pt,
  labelwidth=1.35em,
  labelsep=0.15em,
  itemsep=0.15\baselineskip,
  parsep=0pt,
  topsep=0pt,
  partopsep=0pt,
]
  \item 校级项目，西北大学 IOC 智慧运营中心建设项目，2024.09--2025.06，项目组主要成员。
  \item 校级项目，西北大学导师上岗系统，2024.10--2024.12，项目组主要成员。
\end{enumerate}
\endgroup

% 文末独立空白页（双面印刷时常用；单面下亦保留版式空白）
\newpage
\thispagestyle{empty}
\null
\newpage
